% Generated draft: requires mathematical and source review.
\documentclass[11pt]{article}
\usepackage[margin=25mm]{geometry}
\usepackage{fontspec}
\setmainfont{Latin Modern Roman}
\usepackage{amsmath,amssymb,mathtools,unicode-math}
\setmathfont{Latin Modern Math}
\usepackage{enumitem,needspace,xurl,hyperref}
\setlength{\parindent}{0pt}
\setlength{\parskip}{5pt}
\setlength{\emergencystretch}{3em}
\newcommand{\N}{\mathbb{N}}
\newcommand{\Z}{\mathbb{Z}}
\newcommand{\Q}{\mathbb{Q}}
\newcommand{\R}{\mathbb{R}}
\newcommand{\C}{\mathbb{C}}
\newcommand{\F}{\mathbb{F}}
\newcommand{\A}{\mathbb{A}}
\newcommand{\PP}{\mathbb{P}}
\newcommand{\E}{\mathbb{E}}
\newcommand{\eps}{\varepsilon}
\newcommand{\dd}{\,\mathrm{d}}
\begin{document}
\section*{UnsolvedMath ID: 30002516}
\textbf{Source problem code:} catalogue entry
\subsection{Definitions and mathematical statement}
Let \(\Lambda\) be a noetherian algebra that is finitely generated as a module over its center
\[
Z=Z(\Lambda)=\{z\in\Lambda: z\lambda=\lambda z\text{ for every }\lambda\in\Lambda\}.
\]
The supplied record does not define its cohomology-annihilator ideal or its degree approximations. A conventional candidate formulation is given below; its module-side convention, degree range, and indexing require confirmation from the cited source.

For this candidate formulation, let \(\operatorname{mod}\Lambda\) denote the category of finitely generated right \(\Lambda\)-modules. Each \(\operatorname{Ext}^i_\Lambda(M,N)\) is a \(Z\)-module through the central action. For a \(Z\)-module \(X\), write
\[
\operatorname{ann}_Z(X)=\{z\in Z:zx=0\text{ for every }x\in X\}.
\]
For \(s\in\Z_{\geq 0}\), define the candidate ideals
\[
\operatorname{ca}_s(\Lambda)
=
\bigcap_{\substack{M,N\in\operatorname{mod}\Lambda\\ i\in\Z_{\geq 0},\ i\geq s}}
\operatorname{ann}_Z\!\left(\operatorname{Ext}^i_\Lambda(M,N)\right),
\qquad
\operatorname{ca}(\Lambda)=\bigcup_{s\geq 0}\operatorname{ca}_s(\Lambda).
\]
These definitions describe annihilation in every degree at least \(s\), rather than only in degree \(s\). The supplied phrase ``degree-\(s\) approximation'' does not itself settle this distinction.

Under these conventions, the intended questions are the following, for every algebra satisfying the stated hypotheses:
\begin{enumerate}
\item Does there exist \(c\in\operatorname{ca}(\Lambda)\) that is a nonzerodivisor? If regularity is intended on \(\Lambda\), this means
\[
\text{for every }x\in\Lambda,\qquad cx=0\Longrightarrow x=0.
\]
Since \(c\) is central, left and right multiplication give the same condition. If regularity is intended only in \(Z\), the quantified element \(x\) instead ranges over \(Z\). The supplied record does not specify which requirement is intended.
\item Does the set
\[
\{s\in\Z_{\geq 0}:\operatorname{ca}_s(\Lambda)=\operatorname{ca}(\Lambda)\}
\]
have a least element, and, when it does, what is its value for \(\Lambda\)? The record supplies neither an indexing convention nor a proposed formula or bound for this stabilization index.
\end{enumerate}
The record also leaves unspecified whether ``noetherian'' means left noetherian, right noetherian, or both. Thus the candidate definitions above make the intended questions explicit conditionally; they do not constitute a verified reconstruction of the source formulation.
\subsection{Short English statement}
For a noetherian algebra finite over its center, ask whether its cohomology-annihilator ideal contains a nonzerodivisor and determine the first degree approximation equal to the full ideal. The supplied record leaves the annihilator convention, module side, indexing, and meaning of nonzerodivisor unresolved. Its supplied status is \emph{partially solved}; this draft does not independently verify that status.
\subsection{Sources}
\begin{itemize}
\item Supplied local UnsolvedMath registry record, permanent ID \(30002516\), source problem code \texttt{catalogue entry}, entitled ``Nonzerodivisors and Stabilization of Cohomology Annihilators.'' The recorded proposed year is 2014; no individual proposer is supplied.
\item Supplied report attribution: \emph{Representation Theory of Quivers and Finite Dimensional Algebras} (2014), Oberwolfach Reports, pp.\ 453--529, DOI 10.4171/owr/2014/08. The supplied literature list labels this ``Oberwolfach Report 08/2014'': \url{https://doi.org/10.4171/owr/2014/08}. The background describes these reports as documenting mathematical workshops and their open-problem discussions.
\item Supplied literature reference: ``Iyengar--Takahashi, annihilation and generation'': \url{https://www.math.nagoya-u.ac.jp/~takahashi/ua22.pdf}. No further bibliographic details are supplied.
\item Literature assessment dated exactly \texttt{2026-08-21}: the supplied curation labels the problem \emph{partially solved} and reports that, for broad classes of commutative noetherian and excellent rings, cohomology annihilators and strong generators are related to the singular locus and explicit bounds are proved, while the general nonzerodivisor question and ``least exponent'' remain unsettled. This is supplied dated curation derived from the Oberwolfach material, not an independently verified literature assessment. The background expressly requests independent verification of time-sensitive status and cited claims.
\item The registry supplies the statement label \texttt{corrected\_verified} and describes a manual semantic, typographic, and LaTeX review as confirming a meaning-preserving cleanup. These are supplied provenance claims; they do not resolve the formulation issues identified in this draft.
\end{itemize}
\end{document}
